\documentclass[11pt,a4paper]{article}
\usepackage[T1]{fontenc}
\usepackage[utf8]{inputenc}
\usepackage{lmodern}
\usepackage{amsmath,amssymb,amsthm}
\usepackage[margin=1.15in]{geometry}
\usepackage{longtable,booktabs,array}
\usepackage[hidelinks]{hyperref}
\usepackage{xurl}
\emergencystretch=3em

\newtheorem{theoremx}{Theorem}
\renewcommand{\thetheoremx}{\Alph{theoremx}}
\newtheorem{lemma}{Lemma}[section]
\newtheorem{corollary}[lemma]{Corollary}
\theoremstyle{definition}
\newtheorem{definition}[lemma]{Definition}
\theoremstyle{remark}
\newtheorem{remark}[lemma]{Remark}

\newcommand{\Z}{\mathbb{Z}}
\newcommand{\R}{\mathbb{R}}
\newcommand{\F}{\mathbb{F}}
\newcommand{\Gam}{\Gamma}
\newcommand{\lt}{\ell^2}
\newcommand{\bt}[1]{b^{(2)}_{#1}}
\newcommand{\Hom}{\operatorname{Hom}}
\newcommand{\Dsp}{\mathcal{D}}
\newcommand{\Cay}{\mathcal{C}}

\title{The first $L^2$-Betti number of aspherical manifolds\\ and Poincar\'e duality
groups of dimension at least three\thanks{AI-produced. The method is that of
OpenAI's manuscript \cite{O}, which this paper cites for the main argument and does
not reproduce. The reductions of Section~\ref{sec:reductions} were found by GPT-6
Astra (OpenAI, ChatGPT Pro mode), and the text was written by Claude Fable 5.2
(Anthropic) at effort max; see Section~\ref{sec:provenance} for who proved, wrote
and checked what. No human has checked this text. To the extent that any rights
exist in it, they are dedicated to the public domain under CC0~1.0.}}
\author{}
\date{25 November 2026}

\begin{document}
\maketitle

\begin{abstract}
The $L^2$-Betti numbers $\bt{p}$ of a finite complex measure the square-summable
harmonic $p$-chains on its universal cover. Singer's conjecture says that for a
closed aspherical manifold of dimension $n$, one whose universal cover is
contractible, they vanish outside the middle degree $n/2$. A manuscript of OpenAI dated 25 September 2026,
released on 6 October 2026, proves the conjecture in dimension four; it is not refereed.
Its proof assumes $\bt{1}>0$, builds from functions of finite energy on the
fundamental group $\Gam$ a compact boundary on which $\Gam$ acts, and reaches a
contradiction; dimension four enters only in a preliminary section. We show that
the same argument gives $\bt{1}=\bt{n-1}=0$ for every finite aspherical Poincar\'e
complex of formal dimension $n\ge3$, hence for every closed aspherical topological
manifold of dimension at least three (Theorem~A), and for every finitely presented
Poincar\'e duality group of dimension $n\ge3$ (Theorem~B). To do so we state what
Sections 3--10 of the manuscript prove as a criterion for groups (Theorem~C): a
finitely presented torsion-free group with one end, with
$H^2(\Gam;\F_2\Gam)=0$ and with no splitting over a finite or two-ended subgroup, has
$\bt{1}=0$; and we prove, in every dimension $n\ge3$, the facts about $\Gam$ that the
criterion needs. Theorems A and B depend on the manuscript's argument, which we
cite and do not reproduce. One AI referee has followed that argument line by line
and found no fault; no human has checked it, or this paper.
\end{abstract}

\section{Results}

\subsection{Definitions}\label{sub:defs}
Throughout, $\Gam$ is a discrete group and $k=\F_2$.

\emph{$L^2$-Betti numbers.} Let $C_*$ be a chain complex of left $\Z\Gam$-modules
that are free with a given finite basis in the degrees considered, $C_p=\Z\Gam^{r_p}$.
Completing $\Z\Gam$ to the Hilbert space $\lt(\Gam)$ of square-summable functions makes
each boundary map a bounded operator
$\partial_p:\lt(\Gam)^{r_p}\to\lt(\Gam)^{r_{p-1}}$ that commutes with the action of $\Gam$
on the other side. The space of \emph{harmonic $p$-chains} is
$\mathcal H_p(C)=\ker\partial_p\cap\ker\partial_{p+1}^*$, a closed $\Gam$-invariant
subspace. Its \emph{von Neumann dimension} is
$\sum_j\langle\pi\delta_{1,j},\delta_{1,j}\rangle$, where $\pi$ is the orthogonal
projection onto it and $\delta_{g,j}$ is the unit vector at $g$ in the
$j$-th coordinate; it is a real number $\ge0$, it is $0$ only for the zero space,
and it is the same for two such subspaces $V,W$ between which there is a bounded
$\Gam$-equivariant bijection $T$. Indeed, the dimension equals
$\sum_j\|\pi\delta_{1,j}\|^2$, and $\pi$ commutes with the action of $\Gam$, in
which $g$ moves $\delta_{h,j}$ to $\delta_{gh,j}$, so the dimension is $0$ only if
$\pi=0$. The polar decomposition $T=U|T|$ has a $\Gam$-equivariant unitary
$U:V\to W$. Let $u$ be $U$ on $V$ and zero on the orthogonal complement of $V$.
Then $u^*u$ and $uu^*$ are the projections onto $V$ and $W$, and $u$ and $u^*$
commute with the action, so that
\[
 \sum_j\|u\delta_{1,j}\|^2=\sum_{i,j,g}|\langle u\delta_{1,j},\delta_{g,i}\rangle|^2
 =\sum_{i,j,g}|\langle\delta_{g^{-1},j},u^*\delta_{1,i}\rangle|^2=\sum_i\|u^*\delta_{1,i}\|^2 ;
\]
the two ends are the dimensions of $V$ and of $W$. For the theory of this
dimension see \cite[Chapter~1]{L}. We write
$\bt{p}(C)$ for the dimension of $\mathcal H_p(C)$. The quotient $\ker\partial_p/\overline{\operatorname{im}
\partial_{p+1}}$, the \emph{reduced homology}, is identified with
$\mathcal H_p(C)$ by orthogonal projection.

For a finite connected CW complex $Q$ with fundamental group $\Gam$ and universal
cover $\widetilde Q$, the cellular chain complex $C_*(\widetilde Q)$ is such a
complex, with one lift of each cell as basis, and
$\bt{p}(\widetilde Q)=\bt{p}(C_*(\widetilde Q))$. For a finitely presented group
$\Gam$ we put $\bt{1}(\Gam)=\bt{1}(\widetilde K)$, where $K$ is the two-dimensional
complex of a finite presentation. By Lemma~\ref{lem:comparison} this does not depend
on the presentation, and $\bt{1}(\widetilde Q)=\bt{1}(\Gam)$ for every finite
connected complex $Q$ with fundamental group $\Gam$.

\emph{Poincar\'e complexes.} A \emph{finite Poincar\'e complex of formal dimension
$n$} is a finite connected CW complex $Q$ with a homomorphism $w:\Gam=\pi_1(Q)\to\{\pm1\}$
and a class $[Q]\in H_n(Q;\Z^w)$, where $\Z^w$ is $\Z$ with $\Gam$ acting through $w$,
such that cap product with $[Q]$ is an isomorphism
$H^i(Q;V)\to H_{n-i}(Q;\Z^w\otimes V)$ for every integer $i$ and every left
$\Z\Gam$-module $V$ \cite{W}. The number $n$ need not be the largest dimension of a
cell. The complex is \emph{aspherical} if $\widetilde Q$ is contractible.

\emph{Duality groups.} A group $\Gam$ is a \emph{Poincar\'e duality group of dimension
$n$}, for short a $PD(n)$ group, if the trivial module $\Z$ has a resolution of
finite length by finitely generated projective $\Z\Gam$-modules, $H^i(\Gam;\Z\Gam)=0$ for
$i\ne n$, and $H^n(\Gam;\Z\Gam)$ is infinite cyclic as an abelian group \cite{Da}. Such
a group need not be finitely presented \cite[Theorem~7.15]{Da}. Its $L^2$-Betti
numbers are those of a finite projective resolution (Section~\ref{sub:pd}).

\emph{Ends and splittings.} A finitely generated group has \emph{one end} if its
Cayley graph is infinite and removing any finite set of vertices leaves exactly one
infinite component. A group is \emph{two-ended} if it has an infinite cyclic
subgroup of finite index. A \emph{splitting} of $\Gam$ over a subgroup $J$ is an
action of $\Gam$ on a simplicial tree, without inversions of edges, with one orbit of
edges and no fixed vertex, in which $J$ is the stabilizer of an edge. By Bass--Serre
theory this is the same as a decomposition of $\Gam$ as a non-trivial amalgamated free
product or HNN extension over $J$.

\subsection{Statements}

\begin{theoremx}\label{thm:A}
Let $Q$ be a finite aspherical Poincar\'e complex of formal dimension $n\ge3$,
with any orientation character $w$. Then
\[
 \bt{1}(\widetilde Q)=\bt{n-1}(\widetilde Q)=0 .
\]
Moreover $\bt{0}(\widetilde Q)=\bt{n}(\widetilde Q)=0$ and $\bt{p}(\widetilde Q)=0$
for $p>n$. The same holds for every closed connected aspherical topological
manifold of dimension $n\ge3$.
\end{theoremx}

\begin{theoremx}\label{thm:B}
Let $\Gam$ be a finitely presented Poincar\'e duality group of dimension $n\ge3$.
Then $\bt{1}(\Gam)=\bt{n-1}(\Gam)=0$.
\end{theoremx}

\begin{theoremx}[OpenAI \cite{O}, Sections 3--10]\label{thm:C}
Let $\Gam$ be a finitely presented torsion-free group. Suppose that $\Gam$ has one
end, that $H^2(\Gam;\F_2\Gam)=0$, and that $\Gam$ has no splitting over a finite or
two-ended subgroup. Then $\bt{1}(\Gam)=0$.
\end{theoremx}

For $n=3$ Theorem~\ref{thm:A} gives the vanishing of all $L^2$-Betti numbers, and
for $n=4$ of all but $\bt{2}$, which is the theorem of \cite{O}. For $n\ge5$ it says
nothing about the degrees $2,\dots,n-2$.

\subsection{What is whose}\label{sub:whose}
\begin{itemize}
\item The method, and the whole argument behind Theorem~\ref{thm:C}, are OpenAI's
\cite{O}. That manuscript states its theorem for formal dimension four. Its
Sections 3--10 are written under the assumption $\bt{1}>0$ for a complex of formal
dimension four; Theorem~\ref{thm:C} is what they prove when they are read for their
hypotheses. Section~\ref{sec:criterion} says exactly where they call on dimension
four, and reports the check on which that reading rests.
\item That the method gives every dimension $n\ge3$ was observed in public by
the commenter \texttt{Apprehensive\_Sand951} in the \texttt{r/math} thread
\emph{AI in mathematics, October 10, 2026} \cite{R}. The comment says that,
should the proof of \cite{O} hold up, its method also shows that ``any closed
aspherical manifold of dimension $\ge3$ has no $L^2$-harmonic $1$-forms'', and
that in dimension three this would prove Singer's conjecture without
geometrization. It is conditional on \cite{O} and gives no proof. For a smooth
manifold, the absence of non-zero $L^2$-harmonic $1$-forms on the universal
cover is the equation $\bt{1}=0$, by the $L^2$ Hodge--de Rham theorem of Dodziuk
\cite[Chapter~1]{L}. This paper proves that observation, as
Theorem~\ref{thm:A}; like the observation, the theorem rests on the argument of
\cite{O}.
\item New here are the reductions of Section~\ref{sec:reductions} for every
$n\ge3$ and for $PD(n)$ groups, the statement of Theorem~\ref{thm:C} with the
account of what Sections 3--10 of \cite{O} use, and with them the proofs of Theorems~\ref{thm:A}
and~\ref{thm:B}.
\item Known before: for closed aspherical three-manifolds the vanishing follows
from the computation of Lott and L\"uck \cite[Theorem~0.1]{LL} and geometrization (stated there for
orientable manifolds; a nonorientable closed aspherical three-manifold is
handled by its orientation double cover, since $L^2$-Betti numbers multiply by
the degree of a finite cover);
for right-angled Coxeter groups whose nerve is a flag triangulation of a
generalized homology sphere, Davis and Okun proved $\bt{1}=\bt{n-1}=0$ in every
dimension $n\ge3$ \cite[Theorem~11.3.2]{DO}. In January 2026 the vanishing of
$\bt{1}$ was described as unknown even for closed negatively curved manifolds of
dimension at least four \cite{DHT}.
\end{itemize}

\subsection{Status}
Theorem~\ref{thm:C} is as reliable as the argument of \cite{O}, which is an
unrefereed manuscript. What has been checked, by whom,
and what has not, is in Sections~\ref{sec:provenance} and~\ref{sec:open}. Until a human
reader has gone through Sections 3--10 of \cite{O},
Theorems~\ref{thm:A} and~\ref{thm:B} should be cited with that condition.

\section{The reductions}\label{sec:reductions}

In this section $Q$ is a finite aspherical Poincar\'e complex of formal dimension
$n$ and $\Gam=\pi_1(Q)$. Since $\widetilde Q$ is contractible, $C_*(\widetilde Q)$ is
a resolution of $\Z$ by finitely generated free $\Z\Gam$-modules, and homology and
cohomology of $Q$ with coefficients in a $\Z\Gam$-module are those of $\Gam$.

\begin{lemma}\label{lem:finiteness}
If $n\ge1$, the group $\Gam$ is finitely presented, torsion-free and infinite.
\end{lemma}
\begin{proof}
The two-skeleton of $Q$ gives a finite presentation. Let $m$ be the largest
dimension of a cell of $Q$. If $\Gam$ had an element of prime order $p$, generating
a subgroup $C$, then $C_*(\widetilde Q)$ would be a free resolution of $\Z$ over
$\Z C$ of length $m$, so that $H^j(C;\Z)=0$ for $j>m$; but $H^{2r}(C;\Z)=\Z/p$ for
all $r\ge1$. So $\Gam$ is torsion-free. If $\Gam$ were finite it would be trivial and
$Q$ contractible; duality with $V=\F_2$ would give
$H_n(Q;\F_2)\cong H^0(Q;\F_2)=\F_2$ with $n\ge1$, which is absurd.
\end{proof}

\begin{lemma}\label{lem:duality}
For every left $k\Gam$-module $V$ and every $i$ there is an isomorphism
$H^i(\Gam;V)\cong H_{n-i}(\Gam;V)$. For every subgroup $J$ of $\Gam$ and every $i$,
\[
 H^i(\Gam;k[\Gam/J])\cong H_{n-i}(J;k),
\]
where $k[\Gam/J]$ is the permutation module with basis the left cosets of $J$.
\end{lemma}
\begin{proof}
On a vector space over $\F_2$ the sign $w$ acts trivially, so
$\Z^w\otimes V\cong V$, and the first statement is the duality of $Q$. For the
second, $k[\Gam/J]\cong k\Gam\otimes_{kJ}k$ as left $k\Gam$-modules. If $P_*$ is a
resolution of $k$ by free right $k\Gam$-modules, then
$P_*\otimes_{k\Gam}(k\Gam\otimes_{kJ}k)\cong P_*\otimes_{kJ}k$, and $P_*$ is also a
resolution by free right $kJ$-modules. Taking homology gives
$H_j(\Gam;k[\Gam/J])\cong H_j(J;k)$.
\end{proof}

\begin{corollary}\label{cor:vanishing}
If $n\ge3$, then $H^1(\Gam;k\Gam)=0$ and $H^2(\Gam;k\Gam)=0$, and $H^1(\Gam;k[\Gam/J])=0$ when
$J$ is trivial or infinite cyclic.
\end{corollary}
\begin{proof}
By Lemma~\ref{lem:duality} these groups are $H_{n-1}(1;k)$, $H_{n-2}(1;k)$ and
$H_{n-1}(J;k)$. The first two vanish because $n-1$ and $n-2$ are positive. An
infinite cyclic group has a free resolution of length one, so its homology
vanishes in degrees $\ge2$, and $n-1\ge2$.
\end{proof}

The next two lemmas concern an arbitrary finitely generated group.

\begin{lemma}\label{lem:oneend}
An infinite finitely generated group $\Gam$ with $H^1(\Gam;k\Gam)=0$ has one end.
\end{lemma}
\begin{proof}
Let $G$ be the Cayley graph of a finite generating set, and
$C_1\to C_0\to k\to0$ its augmented chain complex over $k$, which is exact
because $G$ is connected; $C_0$ and $C_1$ are finitely generated free
$k\Gam$-modules. Choose a free module $P_2$ mapping onto the kernel of
$C_1\to C_0$ and continue to a free resolution $P_*$ of $k$ with $P_0=C_0$,
$P_1=C_1$. For $i=0,1$, a $k$-valued cochain $a$ on the $i$-cells of $G$ with
finite support defines the homomorphism
$\Phi(a):C_i\to k\Gam$, $\Phi(a)(c)=\sum_{g\in\Gam}a(g^{-1}c)\,g$; every homomorphism
$C_i\to k\Gam$ is of this form for exactly one $a$ (take the coefficient of the
identity), and $\Phi(da)=\Phi(a)\circ\partial$.

Suppose that for some finite set $A$ of vertices the graph $G-A$ has two infinite
components, and let $f$ be the indicator function of one of them. Then $df$ has
finite support, since an edge on which $f$ changes meets $A$. For a finite cycle
$z$ in $C_1$ and $g\in\Gam$ we have $df(g^{-1}z)=f(\partial g^{-1}z)=0$, so
$\Phi(df)$ vanishes on the image of $P_2$ and is a cocycle of
$\Hom_{k\Gam}(P_*,k\Gam)$. As $H^1(\Gam;k\Gam)=0$, there is a cochain $u$ of finite
support on the vertices with $\Phi(df)=\Phi(u)\circ\partial=\Phi(du)$, that is,
$d(f-u)=0$. Then $f-u$ is constant, although it equals $f$ outside a finite set
and $f$ takes each of its two values infinitely often. So $G-A$ has at most one
infinite component; it has one, because $G$ is infinite and locally finite, so
that $G-A$ has finitely many components.
\end{proof}

\begin{lemma}\label{lem:nosplit}
Let $\Gam$ be a group and $J$ a subgroup with $H^1(\Gam;k[\Gam/J])=0$. Then $\Gam$ has no
splitting over $J$.
\end{lemma}
\begin{proof}
Suppose $\Gam$ acts on a tree $T$ without inversions, with one orbit $E$ of edges,
and that $J$ is the stabilizer of an edge $e$. Then $ge\mapsto gJ$ identifies $E$
with $\Gam/J$ as a $\Gam$-set. Fix a vertex $o$, and for $g\in\Gam$ let
$c(g)\in k[E]$ be the sum of the edges on the geodesic from $o$ to $go$. The path
from $o$ to $go$ followed by the $g$-translate of the path from $o$ to $ho$ ends
at $gho$, and reducing it to the geodesic removes edges in pairs; hence
\[
 c(gh)=c(g)+g\,c(h),
\]
so $c$ is a $1$-cocycle of $\Gam$ with values in $k[E]\cong k[\Gam/J]$. By hypothesis
there is $v\in k[E]$ with $c(g)=gv-v$ for all $g$. A geodesic contains each of its
edges once, so the distance from $o$ to $go$ is the number of edges in $c(g)$,
which is at most twice the number of edges in $v$. The orbit of $o$ is therefore
bounded. A group acting on a tree with a bounded orbit fixes the centre of the
subtree spanned by that orbit, which is a vertex or the midpoint of an edge; as
no edge is inverted, $\Gam$ fixes a vertex. So the action is not a splitting.
\end{proof}

\begin{corollary}\label{cor:group}
If $n\ge3$, the group $\Gam$ is finitely presented and torsion-free, it has one end,
$H^2(\Gam;\F_2\Gam)=0$, and $\Gam$ has no splitting over a finite or two-ended subgroup.
\end{corollary}
\begin{proof}
Lemma~\ref{lem:finiteness}, Corollary~\ref{cor:vanishing} and
Lemma~\ref{lem:oneend} give all but the last statement. A finite subgroup of $\Gam$
is trivial. A two-ended group has a finite normal subgroup with quotient infinite
cyclic or infinite dihedral \cite[Lemma~4.1]{W}; in a torsion-free group the
finite subgroup is trivial and the dihedral case is excluded, so a two-ended
subgroup of $\Gam$ is infinite cyclic. Corollary~\ref{cor:vanishing} and
Lemma~\ref{lem:nosplit} now exclude both kinds of splitting.
\end{proof}

\subsection{The first $L^2$-Betti number and functions of finite energy}

\begin{lemma}\label{lem:comparison}
Let $P_*$ and $P'_*$ be complexes of left $\Z\Gam$-modules, free with given finite
bases in degrees $0,1,2$, each with a surjection $P_0\to\Z$ such that
$P_2\to P_1\to P_0\to\Z\to0$ is exact. Then $\bt{1}(P)=\bt{1}(P')$. If $P_*$ and
$P'_*$ are bounded complexes of finitely generated free modules and are chain
homotopy equivalent, then $\bt{p}(P)=\bt{p}(P')$ for all $p$.
\end{lemma}
\begin{proof}
Since the modules are free and the complexes exact in degrees $0$ and $1$, the
identity of $\Z$ lifts to maps $f_i:P_i\to P'_i$ and $g_i:P'_i\to P_i$,
$i=0,1,2$, commuting with the boundary maps, and there are $h_0:P_0\to P_1$ and
$h_1:P_1\to P_2$ with $g_0f_0-1=\partial h_0$ and
$g_1f_1-1=\partial h_1+h_0\partial$: first $g_0f_0-1$ has image in the kernel of
the augmentation, which is the image of $\partial_1$; then
$g_1f_1-1-h_0\partial$ has image in $\ker\partial_1=\operatorname{im}\partial_2$.
All these maps are matrices over $\Z\Gam$ and extend to bounded $\Gam$-equivariant
operators on the completions. So $f_1$ maps $\ker\partial_1$ to $\ker\partial'_1$
and the closure of $\operatorname{im}\partial_2$ into that of
$\operatorname{im}\partial'_2$, and induces a bounded equivariant map of reduced
homology in degree one; so does $g_1$; and $g_1f_1$ induces the identity, because
$g_1f_1-1$ maps $\ker\partial_1$ into $\operatorname{im}\partial_2$. By symmetry
the two maps are inverse to each other, and the dimensions agree. The second
statement is proved in the same way from a chain homotopy equivalence.
\end{proof}

In particular $\bt{1}(\Gam)$ is well defined, $\bt{1}(\widetilde Q)=\bt{1}(\Gam)$ for
every finite connected CW complex $Q$ with fundamental group $\Gam$, since
$C_*(\widetilde Q)$ is exact in degree one, and $\bt{p}(\widetilde Q)$ is an
invariant of the homotopy type of $Q$, because a homotopy equivalence of finite
complexes lifts to a chain homotopy equivalence of the cellular chain complexes of
the universal covers.

Let now $\Gam$ have a finite presentation with relators of length at most three; any
finite presentation can be brought to this form by adding a generator for each
proper prefix of length at least two of a longer relator. Let $\Cay$ be its Cayley
two-complex, the universal cover of the presentation complex, and $G$ its graph,
with set of edges $E$. For $f:\Gam\to\R$ let $df(e)$ be the difference of the values
of $f$ at the ends of the oriented edge $e$, and put
\[
 \Dsp=\Bigl\{df:\ \sum_{e\in E}|df(e)|^2<\infty\Bigr\},\qquad
 \Dsp_0=\overline{\{df:\ f\text{ of finite support}\}}\subset\lt(E;\R).
\]

\begin{lemma}\label{lem:gradient}
The space $\Dsp$ is closed in $\lt(E;\R)$ and contains $\Dsp_0$, and $\bt{1}(\Gam)>0$ if
and only if $\Dsp\ne\Dsp_0$.
\end{lemma}
\begin{proof}
Write $d_0=\partial_1^*$ and $d_1=\partial_2^*$ for the adjoints of the completed
boundary maps of $C_*(\Cay)$; on real cochains they are the cellular coboundary
maps, and they are bounded because every edge lies on a bounded number of
two-cells. A square-summable real $1$-cochain $a$ with $d_1a=0$ is a cocycle on
the simply connected complex $\Cay$, hence $a=df$ for a function $f$, obtained by
summing $a$ along paths from a base vertex. So $\Dsp=\ker d_1$, which is closed.
Functions of finite support are dense in $\lt(\Gam;\R)$ and $d_0$ is bounded, so
$\Dsp_0$ is the closure of the image of $d_0$. The orthogonal complement of $\Dsp_0$
in $\Dsp$ is therefore $\ker d_1\cap\ker d_0^*=\ker\partial_2^*\cap\ker\partial_1$,
the space of real harmonic $1$-chains. The boundary maps have integer matrices,
so the space of complex harmonic $1$-chains is its complexification, and it is
non-zero exactly when its dimension $\bt{1}(\Gam)$ is positive.
\end{proof}

\subsection{Duality for $L^2$-Betti numbers}

\begin{lemma}\label{lem:l2duality}
Let $Q$ be a finite Poincar\'e complex of formal dimension $n$. Then
$\bt{p}(\widetilde Q)=\bt{n-p}(\widetilde Q)$ for every integer $p$, where
$\bt{p}=0$ for $p<0$. If $\Gam$ is infinite then $\bt{0}(\widetilde Q)=0$.
\end{lemma}
\begin{proof}
Let $C=C_*(\widetilde Q)$, and let $C^{n-*}$ be the complex with
$\Hom_{\Z\Gam}(C_{n-p},\Z\Gam)$ in degree $p$, regarded as a complex of left modules
by means of the anti-automorphism $g\mapsto w(g)g^{-1}$ of $\Z\Gam$. Cap product
with a cycle representing $[Q]$ is defined on chains, by a cellular approximation
of the diagonal, and gives a chain map $C^{n-*}\to C$ of $\Z\Gam$-modules which
induces the duality isomorphisms for $V=\Z\Gam$ \cite[Section~1]{W}; so it is an
isomorphism on homology. Both
complexes are bounded and consist of finitely generated free modules, so the
mapping cone of this map is a bounded exact complex of free modules, hence
contractible, and the map is a chain homotopy equivalence. By
Lemma~\ref{lem:comparison}, $\bt{p}(C^{n-*})=\bt{p}(C)$.

In the dual bases, the boundary map of $C^{n-*}$ in degree $p$ is, up to sign,
the adjoint $\partial_{n-p+1}^*$ conjugated by the unitary operator $U_w$ of
$\lt(\Gam)$ that multiplies the basis vector $\delta_g$ by $w(g)$; this operator
fixes $\delta_1$, so conjugation by it does not change von Neumann dimensions.
Hence the harmonic $p$-chains of $C^{n-*}$ correspond to
$\ker\partial_{n-p+1}^*\cap\ker\partial_{n-p}$, the harmonic $(n-p)$-chains of
$C$, and $\bt{p}(C^{n-*})=\bt{n-p}(C)$.

A harmonic $0$-chain is a square-summable function on the vertices of
$\widetilde Q$ that is killed by $\partial_1^*$, that is, has equal values at the
two ends of every edge. The one-skeleton is connected, so the function is
constant, and it is zero when there are infinitely many vertices.
\end{proof}

\subsection{Poincar\'e duality groups}\label{sub:pd}
Let $\Gam$ be a $PD(n)$ group, $P_*$ a resolution of $\Z$ of finite length by finitely
generated projective left $\Z\Gam$-modules, and $w:\Gam\to\{\pm1\}$ the character
through which $\Gam$ acts on the right on $H^n(\Gam;\Z\Gam)\cong\Z$. A finitely generated
projective module is the image $\Z\Gam^r e$ of an idempotent matrix $e$; its
completion is the range of $e$ on $\lt(\Gam)^r$, a closed invariant subspace, and
the definitions of Section~\ref{sub:defs} and the second statement of
Lemma~\ref{lem:comparison} apply to complexes
of such modules: a homomorphism between two of them is given by a matrix over
$\Z\Gam$, and the lifts and homotopies in the proof of that lemma exist for
projective modules by their lifting property. We put
$\bt{p}(\Gam)=\bt{p}(P)$; for $p=1$ and $\Gam$ finitely presented this is the number
defined before, by the argument of Lemma~\ref{lem:comparison}.

\begin{lemma}\label{lem:pd}
Let $\Gam$ be a $PD(n)$ group with $n\ge1$. Then $\Gam$ is torsion-free and infinite;
$H^i(\Gam;V)\cong H_{n-i}(\Gam;V)$ for every $k\Gam$-module $V$ and every $i$, so that
the conclusion of Lemma~\ref{lem:duality} holds for $\Gam$; and
$\bt{p}(\Gam)=\bt{n-p}(\Gam)$ for every $p$, with $\bt{0}(\Gam)=0$.
\end{lemma}
\begin{proof}
Torsion is excluded as in Lemma~\ref{lem:finiteness}, since $P_*$ restricts to a
projective resolution of finite length over a subgroup of prime order. The
trivial group has $H^0(\Gam;\Z\Gam)=\Z$, so it is not a $PD(n)$ group for $n\ge1$.

Let $D_*$ be the complex of right modules with $\Hom_{\Z\Gam}(P_{n-j},\Z\Gam)$ in
degree $j$. Its homology in degree $j$ is $H^{n-j}(\Gam;\Z\Gam)$, which is $\Z^w$ for
$j=0$ and zero otherwise. If $D_j\ne0$ for some $j<0$, the boundary map onto the
lowest such module is surjective and splits, and a contractible summand can be
removed; repeating this, $D_*$ is chain homotopy equivalent to a resolution of the
right module $\Z^w$ by finitely generated projective modules. The complex
$\overline P^{\,w}_*$, which is $P_*$ with the right action $x\cdot g=w(g)g^{-1}x$,
is another such resolution. So $D_*$ and $\overline P^{\,w}_*$ are chain homotopy
equivalent.

Tensoring with a left $k\Gam$-module $V$: for finitely generated projective $L$ the
natural map $\Hom_{\Z\Gam}(L,\Z\Gam)\otimes_{\Z\Gam}V\to\Hom_{\Z\Gam}(L,V)$ is an
isomorphism, so the homology of $D_*\otimes_{\Z\Gam}V$ in degree $j$ is
$H^{n-j}(\Gam;V)$; that of $\overline P^{\,w}_*\otimes_{\Z\Gam}V$ is
$H_j(\Gam;\Z^w\otimes V)=H_j(\Gam;V)$. This is the duality.

For the $L^2$-Betti numbers, the chain homotopy equivalence and
Lemma~\ref{lem:comparison} give $\bt{j}(D)=\bt{j}(\overline P^{\,w})=\bt{j}(\Gam)$,
the last equality as in the proof of Lemma~\ref{lem:l2duality}. For a projective
module with idempotent $e$ and completion $H=\operatorname{im}e$, the completion
of the dual module is $\operatorname{im}e^*$, and orthogonal projection onto $H$
maps it isomorphically to $H$, with inverse the restriction of $e^*$; under these
isomorphisms the boundary maps of $D_*$ become the adjoints of those of $P_*$. So
$\bt{j}(D)=\bt{n-j}(\Gam)$. 

Finally $\bt{0}(\Gam)=0$. A harmonic $0$-chain of $P_*$ is an element $\xi$ of
the completion of $P_0$ that is orthogonal to the image of $P_1$, which is the
kernel of the augmentation $\varepsilon:P_0\to\Z$. Choose $x\in P_0$ with
$\varepsilon(x)=1$. For $g\in\Gam$ the element $gx-x$ lies in that kernel, so
$\langle\xi,gx\rangle=\langle\xi,x\rangle$. As $x$ has finite support and $\xi$
is square-summable, $\langle\xi,gx\rangle$ tends to $0$ along every sequence of
distinct elements $g$ of the infinite group $\Gam$. Hence
$\langle\xi,x\rangle=0$, so $\xi$ is orthogonal to $P_0=\ker\varepsilon+\Z x$,
which spans a dense subspace of the completion, and $\xi=0$.
\end{proof}

\section{The criterion}\label{sec:criterion}

\subsection{The manuscript}
The theorem of \cite{O} is: \emph{for a finite aspherical Poincar\'e complex of
formal dimension four, with any orientation character, $\bt{p}(\widetilde Q)=0$
for all $p\ne2$.} Its Section~2 proves seven lemmas about $\Gam=\pi_1(Q)$ in
dimension four, numbered 2.1 to 2.7 there: finiteness properties of $\Gam$ and a four-dimensional model of $Q$; the
comparison of resolutions; duality for $L^2$-Betti numbers; the vanishing of
$H^i(\Gam;\F_2\Gam)$ for $i=1,2$; one end; the absence of splittings; and the
passage from $\bt{1}>0$ to $\Dsp\ne\Dsp_0$. Our Section~\ref{sec:reductions} replaces these for all $n\ge3$.

Sections 3 to 10 of \cite{O} assume $\bt{1}(\widetilde Q)>0$ and reach a
contradiction through 38 numbered statements, 3.1 to 10.3. In outline: from a
function of finite energy that is not constant at infinity they build a compact
metrizable space $X=\Gam\sqcup B$ on which $\Gam$ acts (Section~3); the action on the
boundary $B$ is a minimal convergence action (Section~4); $B$ is a continuum of
covering dimension at most one with vanishing first \v{C}ech cohomology over
$\F_2$ (Section~5); if $B$ has cut points, they are fixed points of parabolic
elements and are the vertices of a tree, and a block $L$ of $B$ without cut points
inherits these properties (Sections 6 and 7); $L$ is locally connected (Sections 8
and 9); and a locally connected continuum of dimension at most one with vanishing
first \v{C}ech cohomology has a cut point (Section~10).

\subsection{Where Sections 3--10 call on Section 2}\label{sub:interface}
Consider the following statements about a group $\Gam$ with a finite presentation
whose relators have length at most three, with Cayley two-complex $\Cay$.
\begin{itemize}
\item[(G0)] $\Dsp\ne\Dsp_0$. This is the assumption to be refuted.
\item[(G1)] $\Gam$ is finitely presented: $\Cay$ is simply connected and each edge
lies on a bounded number of two-cells.
\item[(G2)] $\Gam$ has one end.
\item[(G3)] $\Gam$ is torsion-free.
\item[(G4)] $H^2(\Gam;\F_2\Gam)=0$.
\item[(G5)] $\Gam$ has no splitting over a finite or two-ended subgroup.
\end{itemize}
Sections 3--10 of \cite{O} refer to its Section~2 in five places, and nowhere else
do they use the complex $Q$, its dimension, or duality:
\begin{enumerate}
\item the opening of Section~3 lists its inputs: the group, the complex $\Cay$ and
its graph, with (G1) and (G2), and the strict inclusion (G0);
\item the proof of Proposition~4.5 uses (G2) and (G3);
\item the proof of Proposition~5.3 cites Lemma~2.4 of \cite{O}, twice, and only
for (G4): a $2$-cochain of finite support of the form $d\alpha$ vanishes on the
two-cycles of $\Cay$, hence is a cocycle on a free resolution that extends the
chains of $\Cay$, and by (G4) it is the coboundary of a cochain of finite support;
\item the opening of Section~6 assumes (G1), (G2), (G3) and (G5), and these are
what Sections 6 and 7 use;
\item the last proof, after statement 10.3, uses Lemma 2.3 of \cite{O}, duality
for $L^2$-Betti numbers, to pass from $\bt{1}=0$ to the other degrees.
\end{enumerate}
Read with (G0)--(G5) as hypotheses in place of items 1--4, Sections 3--10 of
\cite{O} prove that (G1)--(G5) contradict (G0). With Lemma~\ref{lem:gradient},
this is Theorem~\ref{thm:C}.

Table~\ref{tab:audit} lists the 38 statements with the hypotheses each proof
uses directly, besides earlier statements of the table. The statements are
described in our words; for their exact form see \cite{O}.

\begin{longtable}{@{}p{0.07\textwidth}p{0.62\textwidth}>{\raggedright\arraybackslash}p{0.23\textwidth}@{}}
\caption{The numbered statements of Sections 3--10 of \cite{O} and what their
proofs use.}\label{tab:audit}\\
\toprule
No. & Statement, in brief & Uses directly\\
\midrule
\endfirsthead
\toprule
No. & Statement, in brief & Uses directly\\
\midrule
\endhead
\bottomrule
\endfoot
3.1 & a function $f$ of finite energy with values in $[0,1]$, not constant at infinity, whose level sets have connected sides and connected cuts & (G0), (G1)\\
3.2 & cuts of the coordinate functions have finitely many components & 3.1\\
3.3 & boxes with generic sides have finitely many graph components; a countable invariant basis of such boxes & (G1), 3.2\\
3.4 & the localized coordinates have finite energy, and their cuts lie in cuts of the first coordinates & 3.3\\
3.5 & the boundary $B$ of $X=\Gam\sqcup B$ is a continuum with more than one point & (G2), 3.1\\
3.6 & vertices near a point of $B$ are joined by graph paths near that point & 3.3, 3.4\\
4.1 & the total length of a translated cut tends to zero; a bound for the diameter of a connected cut & (G1)\\
4.2 & translated cuts collapse to the limit of the translating elements & 3.2, 3.4, 4.1\\
4.3 & if $g_j\to a$ and $g_j^{-1}\to b$ then $g_j\to a$ uniformly on compact subsets of $X\setminus\{b\}$ & 4.2\\
4.4 & for the fixed point $p$ of a parabolic element, $g$ and $gp$ become close as $g$ leaves finite sets & 4.3\\
4.5 & every orbit in $B$ is dense & (G2), (G3), 3.5, 4.3, 4.4, \cite{BoT}\\
5.1 & almost every cut accumulates on a subset of $B$ of one-dimensional Hausdorff measure zero & 3.2, 3.4, 4.1\\
5.2 & $B$ has covering dimension at most one & 3.6, 5.1\\
5.3 & $\check H^1(B;\F_2)=0$ & (G1), (G4), 3.6\\
6.1 & Bowditch's completed pretree of cut points and the dendritic quotient of $B$ & \cite{BoT}\\
6.2 & the dendritic quotient is a point, and cut points have infinite stabilizers without loxodromic elements & (G1), (G2), (G3), (G5), \cite{BoC}, \cite{BoT}\\
6.3 & the fixed points of a parabolic element in the pretree & 6.1, \cite{BoT}\\
6.4 & an invariant class of the pretree with finite intervals & (G1), (G2), (G3), 6.1--6.3, \cite{BoT}\\
6.5 & a simplicial tree with finite quotient containing the cut points, with infinite edge stabilizers & (G1), (G2), (G5), 6.1, 6.2, 6.4\\
7.1 & a graph on the cut points next to a vertex of the tree that is not a cut point & (G1), 6.5\\
7.2 & paths in that graph near a point of $B$ & 3.6, 4.4, 7.1\\
7.3 & a retraction of $B$ onto a block $L$; $L$ has no cut point, dimension at most one, $\check H^1(L;\F_2)=0$ & 4.5, 5.2, 5.3, 6.5, 7.1, 7.2\\
7.4 & the common setting: a subgroup $H$, the space $H\sqcup L$, a graph with finitely many orbits & (G3), 3.6, 4.3, 4.4, 6.2, 7.1--7.3\\
8.1 & a locally constant function to $\F_2$ on a punctured neighbourhood in $L$ extends over the point & properties of $L$ in 7.3\\
8.2 & points of $L$ away from deleted vertices have a well-defined graph component & 7.4, 8.1\\
8.3 & an invariant system of annuli & 7.4\\
8.4 & the annular depth & 8.3\\
9.1 & translates of a closed set by elements of large depth shrink to a point & 8.3, 8.4\\
9.2 & a comparison of depths, uniform away from that point & 8.3, 8.4, 9.1\\
9.3 & component labels at each depth & 7.4, 8.2, 8.3, 8.4\\
9.4 & few deleted vertices after translation, away from that point & 7.4, 9.1, 9.2\\
9.5 & labels of two points agree eventually & 7.3, 8.2, 9.4\\
9.6 & a threshold for the majority label of three points & 7.4, 9.1, 9.5\\
9.7 & beyond the threshold, greater depth does not separate equal labels & 7.4, 9.3, 9.6\\
9.8 & $L$ is locally connected & 7.4, 9.3, 9.7\\
10.1 & restriction is onto in first \v{C}ech cohomology when the dimension is at most one & ---\\
10.2 & connected open subsets of a locally connected compact metric space are path connected & ---\\
10.3 & a locally connected continuum with more than one point, dimension at most one and $\check H^1=0$ has a cut point & 10.1, 10.2, \cite{K}\\
\end{longtable}

\subsection{The outside results used by the manuscript}
Besides elementary analysis and topology, Sections 3--10 of \cite{O} use the
following published results.
\begin{itemize}
\item Bowditch \cite[Lemma~1.4]{BoC}: if a finitely presented one-ended group with no
infinite torsion subgroup acts as a minimal convergence group on a dendrite,
in such a way that every two-ended subgroup over which the group splits fixes a
point of the dendrite, then the dendrite is a single point. Under (G5) the condition on
splittings is empty. This is the only place where splittings over two-ended
subgroups enter.
\item Bowditch \cite{BoT}: Theorem~6.1, that for a minimal convergence action of a
finitely generated one-ended group on a
continuum a cut point whose stabilizer is finite or contains a loxodromic
element makes the dendritic quotient non-trivial; Lemma~6.2, that in a
convergence action on a compact Hausdorff space with at least three points an
element of infinite order is parabolic, with one fixed point, or loxodromic,
with two, and Proposition~6.4 with the paragraph before it, that a point
stabilizer containing a loxodromic element is virtually cyclic (both are quoted
after statement 4.3 and used from statement 4.5 on); and, for statement 6.1 of
the table, Theorem~3.19, Lemmas 3.13, 3.28, 3.34,
4.2--4.4, Theorem~5.23, and Lemmas 6.5, 6.6, 6.8--6.11 and 6.13.
\item Moore's theorem that a continuous image of an interval contains an arc
between any two of its points, in the form of \cite{K}.
\end{itemize}

\subsection{The check}\label{sub:check}
The reading in Section~\ref{sub:interface} rests on the following checks;
details and dates are in Section~\ref{sec:provenance}. The text read line by
line was a reproduction of Sections 3--10 of \cite{O}, compared with the source
files of \cite{O}. Sections 3, 7 and 8 are identical. In Section~4 the
reproduction adds one sentence; in Section~5 it spells out the last step of the
proof of 5.3; in Section~9 it rewrites the limit argument at the end of 9.8; in
Section~10 it gives statement 10.1 for the coefficients $\F_2$ only and writes
the last proof for general $n$; the other differences, in Sections 6, 9 and 10,
are of wording. One AI referee followed the proof of
every one of the 38 statements to the end, those of Sections 3, 4, 5, 8 and 9
twice, with attention to the five places of Section~\ref{sub:interface} and to
every use of (G1)--(G5), and found no step that fails and no further use of
Section~2. An earlier AI referee had passed Sections 3, 4, 5, 7, 8 and 10
and the arguments of Section~6, without opening the results of Bowditch, and had
examined Section~9 only in part. A second reader, Claude Opus~5.5 at effort
high, then read Sections 8 and 9 line by line in the manuscript's own source
before seeing the first referee's report, every case of every proof, and found
that all twelve statements hold; it agrees with the first referee on each of
them, and finds both forms of the limit argument at the end of 9.8 correct. A third AI reader compared
Section~\ref{sub:interface} and Table~\ref{tab:audit} with the manuscript
itself, its PDF and its source: the numbers and the descriptions of the
statements, the five places, and, statement by statement, the references made
in each proof. The results of Bowditch and of Kamiya were opened at the
statements listed above and are used as stated. Their proofs were not examined,
except that of Lemma~1.4 of \cite{BoC}, which the third reader read; what it
rests on is in Section~\ref{sec:open}.

\section{Proofs of Theorems A and B}

\begin{proof}[Proof of Theorem~\ref{thm:A}]
By Corollary~\ref{cor:group} the group $\Gam=\pi_1(Q)$ satisfies the hypotheses of
Theorem~\ref{thm:C}, so $\bt{1}(\Gam)=0$. By Lemma~\ref{lem:comparison},
$\bt{1}(\widetilde Q)=\bt{1}(\Gam)=0$. Lemma~\ref{lem:l2duality} gives
$\bt{n-1}(\widetilde Q)=\bt{1}(\widetilde Q)=0$ and
$\bt{n}(\widetilde Q)=\bt{0}(\widetilde Q)=0$, since $\Gam$ is infinite by
Lemma~\ref{lem:finiteness}, and $\bt{p}(\widetilde Q)=\bt{n-p}(\widetilde Q)=0$ for
$p>n$.

A closed connected topological manifold $M$ of dimension $n$ has the homotopy
type of a finite CW complex: it is a compact absolute neighbourhood retract
\cite[Theorems 3.2 and 3.3]{Ha}, and such a space is homotopy equivalent to a
finite complex \cite[Corollary~5.3]{We}. It satisfies Poincar\'e duality with its
orientation character and all local coefficients
\cite[Theorem~10.2 and Remark~10.3]{Sp}: the duality isomorphism there is
defined from the Thom class, for \v{C}ech--Alexander cohomology; for the pair
$(M,\emptyset)$ it corresponds to cap product with the fundamental class, and
on a manifold \v{C}ech--Alexander and singular cohomology agree. This passes to
the finite complex.
If $M$ is aspherical, so is the complex, and the first part applies; the
$L^2$-Betti numbers of $M$ are those of the complex, which by
Lemma~\ref{lem:comparison} do not depend on its choice.
\end{proof}

\begin{proof}[Proof of Theorem~\ref{thm:B}]
By Lemma~\ref{lem:pd}, $\Gam$ is torsion-free and infinite and satisfies the
conclusion of Lemma~\ref{lem:duality}, hence that of
Corollary~\ref{cor:vanishing}, whose proof uses nothing else. Lemmas
\ref{lem:oneend} and \ref{lem:nosplit} then show, as in
Corollary~\ref{cor:group}, that $\Gam$ has one end and no splitting over a finite
or two-ended subgroup, and $H^2(\Gam;\F_2\Gam)=0$. Since $\Gam$ is finitely presented
by hypothesis, Theorem~\ref{thm:C} gives $\bt{1}(\Gam)=0$, and Lemma~\ref{lem:pd}
gives $\bt{n-1}(\Gam)=\bt{1}(\Gam)=0$.
\end{proof}

\section{Remarks}

\begin{remark}[Dimension two]
The bound $n\ge3$ is sharp. For a closed orientable surface of genus $g\ge2$,
$H^2(\Gam;\F_2\Gam)\cong H_0(1;\F_2)=\F_2$ by Lemma~\ref{lem:duality}, so (G4) fails,
and indeed $\bt{1}=2g-2$. In the proofs above the dimension is used in
Corollary~\ref{cor:vanishing} only, through $n-2>0$ and $n-1\ge2$.
\end{remark}

\begin{remark}[Finite presentation]
Theorem~\ref{thm:B} assumes that the group is finitely presented. The argument of
\cite{O} uses a finite presentation from its first step, in the connectedness of
cuts, and Bowditch's lemma assumes one.
\end{remark}

\begin{remark}[What Sections 3 and 4 use]
Statements 3.1 to 4.5 of Table~\ref{tab:audit} use (G0)--(G3) only: neither the
vanishing of $H^2(\Gam;\F_2\Gam)$ nor the absence of splittings. What they say about
arbitrary finitely presented torsion-free groups with one end and $\bt{1}>0$,
which exist, is left to a separate note.
\end{remark}

\section{Provenance and checks}\label{sec:provenance}

This paper is AI-produced. No human has proved, written or checked any part of
it.

\paragraph{Sources.} The argument for Theorem~\ref{thm:C} is that of the
manuscript \cite{O} of OpenAI, dated 25 September 2026 and released on 6 October
2026; we cite it and do not reproduce it. That the method gives every dimension
$n\ge3$ was observed in the comment \cite{R}, which gives no proof and is
earlier than the run described next. The reductions of Section~2 of \cite{O} in
every dimension $n\ge3$, with their proofs, the group criterion in the form of
Theorem~\ref{thm:C}, and the extension to finitely presented $PD(n)$ groups are
due to a run of GPT-6 Astra (OpenAI, ChatGPT Pro mode) of 22--23 October 2026, in
answer to a question we had put to it after reading \cite{O}; that run also
reproduced Sections 3--10 of \cite{O} with a table of the inputs of each
statement.

\paragraph{The text.} The present text was written by Claude Fable 5.2 at effort
max. Section~\ref{sec:reductions} follows the mathematics of the run and is
written afresh; it defines $L^2$-Betti numbers through harmonic chains, where the
run used modules over the von Neumann algebra, and its Lemma~\ref{lem:comparison}
is proved accordingly. Section~\ref{sec:criterion}, with
Table~\ref{tab:audit}, is the writer's reading of \cite{O}.

\paragraph{Checks.} The run's version of the reductions was checked in two
passes: its Lemmas on finiteness, duality for $L^2$-Betti numbers, duality over
$\F_2$, ends and splittings by Claude Opus 5.5 at effort high (24 November 2026),
and its comparison lemma, its gradient lemma and its treatment of $PD(n)$ groups
by Claude Fable 5.2 at effort max (25 November 2026). Sections 3--10 of \cite{O},
in the run's reproduction, were read by Claude Opus 5.5 at effort max (Message
Batches API, 23 October 2026), which passed Sections 3, 4, 5, 7, 8 and 10
and the arguments of Section~6, left the citations of Bowditch unopened and
examined Section~9 only in part; and by Claude
Fable 5.2 at effort max (25 November 2026), as described in
Section~\ref{sub:check}. The second of these readers is the writer of this text.
The finished text was then checked independently by Claude Opus 5.5 at effort
high (25 November 2026). It followed every proof of
Section~\ref{sec:reductions} to the end and found no error. It compared
Section~\ref{sec:criterion} with the manuscript's PDF and source; this
confirmed the numbering, the descriptions of the statements and the five places
of Section~\ref{sub:interface}, and it corrected the inputs in five rows of
Table~\ref{tab:audit} (4.5, 6.3, 6.5, 7.4 and 9.6) and the list of outside
results, where Lemma~6.2 and Proposition~6.4 of \cite{BoT} were missing. It read
the proof of Lemma~1.4 of \cite{BoC}, and it opened \cite{W}, \cite{LL},
\cite{DO}, \cite{Da}, \cite{DHT}, \cite{Ha}, \cite{We} and \cite{Sp} at the
statements used. Its repairs were made, some of them in a different form; none changes a theorem
or a lemma.
Passages written or reworded after that check were read by the writer and
then, with every page of this text, by Claude Opus~5.5 at effort max (25
November 2026); the mathematical ones are the proof in Section~\ref{sub:defs}
that the dimension is the same for $V$ and $W$, and the last paragraph of the
proof of Lemma~\ref{lem:pd}. That reading also removed from the description of
Theorem~6.1 of \cite{BoT} in Section~\ref{sec:criterion} a hypothesis on
torsion which Bowditch's statement does not contain. Nothing
is formalized. The record of every check, with model, effort and date, is in the
file \texttt{CHECK\_EVENTS.tsv} that accompanies this paper in its workbench,
reached from \url{https://github.com/KokunoYumeto/mathematics-workbenches};
\texttt{PROVENANCE.md} lists every contributor and says in which copy each
citation was confirmed, and which citations have not been opened by a reader of
ours.

\section{What remains open}\label{sec:open}
\begin{enumerate}
\item A reading of the whole of Sections 3--10 of \cite{O} by a human who knows
convergence groups.
\item The results on which Bowditch's Lemma~1.4 of \cite{BoC} rests. Its proof
has been read by one AI reader. It uses the construction of an $\R$-tree from a
minimal convergence action on a dendrite (Levitt; Bowditch) and a theorem of
Tukia on free abelian subgroups of convergence groups, which no reader of ours
has opened, and the theorem of Bestvina and Feighn on stable actions of
finitely presented groups on $\R$-trees, which has been opened at its
statement.
\item The degrees $2,\dots,n-2$ for $n\ge5$, on which the method gives nothing.
\item Whether Theorem~\ref{thm:B} holds for $PD(n)$ groups that are not finitely
presented.
\end{enumerate}

\begin{thebibliography}{99}
\bibitem{BoC} B.~H.~Bowditch, \emph{Connectedness properties of limit sets},
Trans.\ Amer.\ Math.\ Soc.\ \textbf{351} (1999), 3673--3686; Lemma~1.4 (author's
manuscript, p.~6).
\bibitem{BoT} B.~H.~Bowditch, \emph{Treelike structures arising from continua and
convergence groups}, Mem.\ Amer.\ Math.\ Soc.\ \textbf{139} (1999), no.~662;
numbering of the author's manuscript: Lemma~3.13, Theorem~3.19, Lemmas 3.28, 3.34,
4.2--4.4, 5.15, Theorem~5.23, Theorem~6.1, Lemma~6.2, Proposition~6.4,
Lemmas 6.5, 6.6, 6.8--6.11, 6.13.
\bibitem{Da} M.~W.~Davis, \emph{Poincar\'e duality groups}, in: Surveys on Surgery
Theory, vol.~1, Ann.\ of Math.\ Stud.\ \textbf{145}, Princeton Univ.\ Press, 2000,
167--193; Definitions 3.1 and 3.2, Theorems 3.3 and 7.15.
\bibitem{DO} M.~W.~Davis and B.~Okun, \emph{Vanishing theorems and conjectures for
the $\ell^2$-homology of right-angled Coxeter groups}, Geom.\ Topol.\ \textbf{5}
(2001), 7--74; Theorem~11.3.2.
\bibitem{DHT} L.~F.~Di Cerbo, H.~Hunter and A.~K.~Thrasher, \emph{Price inequality
and the growth of harmonic functions on non-positively curved manifolds},
arXiv:2601.08804; the introduction.
\bibitem{Ha} O.~Hanner, \emph{Some theorems on absolute neighborhood retracts},
Ark.\ Mat.\ \textbf{1} (1951), 389--408; Theorems 3.2 and 3.3.
\bibitem{K} H.~Kamiya, \emph{An elementary proof of a theorem of L.~Moore's on
continuous curves}, Japan.\ J.\ Math.\ \textbf{7} (1930), 301--304; the theorem on
p.~301.
\bibitem{LL} J.~Lott and W.~L\"uck, \emph{$L^2$-topological invariants of
$3$-manifolds}, Invent.\ Math.\ \textbf{120} (1995), 15--60; Theorem~0.1.
\bibitem{L} W.~L\"uck, \emph{$L^2$-Invariants: Theory and Applications to Geometry
and $K$-Theory}, Ergeb.\ Math.\ Grenzgeb.\ (3) \textbf{44}, Springer, 2002; Chapter~1.
\bibitem{O} OpenAI, \emph{The Singer conjecture in dimension four}, manuscript
dated 25 September 2026, with its \TeX\ source, at
\url{https://github.com/openai/math/tree/main/preprints/The-Singer-conjecture-in-dimension-four-September-25-2026}
(Apache License 2.0; released there on 6 October 2026); Sections 2--10.
\bibitem{R} \texttt{Apprehensive\_Sand951}, comment \texttt{pf4k77i} in the thread
\emph{AI in mathematics, October 10, 2026}, \texttt{r/math},
\url{https://www.reddit.com/r/math/comments/1x2iy00/comment/pf4k77i/}.
\bibitem{Sp} E.~Spanier, \emph{Singular homology and cohomology with local
coefficients and duality for manifolds}, Pacific J.\ Math.\ \textbf{160} (1993),
165--200; Theorem~10.2 and Remark~10.3.
\bibitem{W} C.~T.~C.~Wall, \emph{Poincar\'e complexes.~I}, Ann.\ of Math.\ (2)
\textbf{86} (1967), 213--245; Section~1 and Lemma~4.1.
\bibitem{We} J.~E.~West, \emph{Mapping Hilbert cube manifolds to ANR's: a solution
of a conjecture of Borsuk}, Ann.\ of Math.\ (2) \textbf{106} (1977), 1--18;
Corollary~5.3.
\end{thebibliography}

\end{document}
